
AI alignment research evaluates whether AI systems conform to human values and preferences, measured through metrics such as helpfulness ratings and harmlessness scores. This evaluation paradigm measures AI performance without considering reciprocal effects on human behavior during collaboration. In high-stakes deployment contexts such as medical diagnosis and code review, human oversight quality functions as a safety layer, yet current evaluation frameworks do not measure whether alignment interventions affect this oversight.

Human factors research has documented automation bias—the tendency for humans to over-rely on automated systems, reducing verification effort and error detection sensitivity. This phenomenon has been observed in aviation, medical diagnosis, and process control. A proposed causal mechanism suggests that increased AI compliance may reduce observable disagreement signals, which could increase perceived reliability beyond objective accuracy, leading to reduced human verification effort through Bayesian updating, potentially degrading error detection sensitivity over time.

This causal chain—policy-layer compliance manipulation → reduced disagreement → trust miscalibration → verification effort reduction → oversight degradation—suggests that alignment evaluation should measure bidirectional effects. However, testing this mechanism requires addressing three challenges: (1) compliance modulation must not confound capability changes, requiring architectural separation; (2) oversight quality must be quantified without confounding task difficulty or AI capability; (3) temporal dynamics must be tracked across timescales from hours to months.

Constitutional AI's policy-layer architecture provides a candidate approach for architectural separation. System prompts or constitutional rules modulate compliance strength while base model weights remain frozen. This architectural approach enables testable capability invariance via statistical validation (ICC > 0.95, ANOVA p > 0.05, Cohen's f < 0.10), though this is an unverified prerequisite assumption requiring empirical testing. Human oversight quality can be measured via signal detection d' on seeded within-distribution errors, isolating error detection sensitivity from task accuracy.

\textbf{Contributions.} We contribute:

\begin{enumerate}
\item \textbf{Operationalization of bidirectional alignment dynamics} as measurable variables: AI Compliance Elasticity (ACE) via policy-layer parameter λ ∈ {0.2, 0.4, 0.6, 0.8, 1.0} and Human Oversight Retention (HOR) via signal detection d' on seeded within-distribution errors.
\end{enumerate}

\begin{enumerate}
\item \textbf{Dependency-driven verification protocol} decomposing the main hypothesis into five sub-hypotheses (h-e1 → h-m1 → h-m2 → h-m3 → h-m4) with pre-specified success criteria, failure routing decisions, and prerequisite dependencies.
\end{enumerate}

\begin{enumerate}
\item \textbf{Infrastructure implementation} with MMLU (57 subjects, 14,042 questions from HuggingFace cais/mmlu) and HumanEval (164 problems from OpenAI human-eval), automated quality control (mock data detection and correction with 6/6 verification checks), and statistical analysis pipeline (ICC/ANOVA/Cohen's f, signal detection d').
\end{enumerate}

\textbf{Status.} This work presents methodological infrastructure without empirical validation. No experiments were executed due to API resource constraints (estimated $1,620 for full MMLU + HumanEval evaluation across 5 compliance levels, approximately 4 hours runtime). All five predictions (P1: capability invariance, P2: ACE-HOR coupling functional form, P3: perceived reliability mediation, P4: metacognitive intervention effectiveness, P5: equilibrium region existence) remain inconclusive. All five key assumptions remain unverified. The contribution is demonstrating how to measure bidirectional alignment dynamics, not what those dynamics are.
